\section{Actual total-space section norms and weighted original Laurent coefficients}

\begin{definition}\label{def:native-tc-weighted-laurent-section-norm}
\lean{EnvelopeProperty.NativeTestConfiguration.TestConfiguration.originalPowerSectionNorm,EnvelopeProperty.NativeTestConfiguration.TestConfiguration.weightedLaurentPowerSectionNorm}
\leanok
Let $M$ be the original linearized invertible module on a test
configuration and let $L$ be its original invertible fibre-one module.
For every original total-space section $s$ of $M^{\otimes n}$, let
$s_i\in H^0(X,L^{\otimes n})$ be its actual Laurent coefficients
after the whole original product transport. For an original trivially
valued point $x$ and $r>0$, define
\[
 W_{n,x,r}(s)=\max_i\{\lVert s_i(x)\rVert_{L^{\otimes n}}r^i\}.
\]
The maximum is over the actual finite support; its empty value is zero.
Here the norms on the original proper $X$ are its genuine model norms.
\uses{thm:native-tc-global-laurent-sections}
\end{definition}
\begin{proof}
Original test-configuration properness gives properness of $X$.
The whole Laurent coefficient equivalence gives an actual finitely
supported family of original sections in every degree, including zero.
Every tensor power of the invertible original $L$ is invertible,
so its true proper-model norm is defined. The trivial base valuation
is bounded by one. Thus the indicated finite maximum is defined for
every original section and every positive radius, retaining negative
integer exponents and zero sections.
\end{proof}

\begin{theorem}\label{thm:native-tc-weighted-laurent-true-section-norms}
\lean{EnvelopeProperty.NativeTestConfiguration.TestConfiguration.exists_weightedLaurentPowerSectionNorm_factor}
\leanok
For every $n\geq0$, $x$ and $0<r\leq1$, there is one number
$c_{n,x,r}>0$ such that, simultaneously for every original total-space
section $s$ of $M^{\otimes n}$,
\[
 \lVert s\rVert_{\mathcal X,r,x}=c_{n,x,r}W_{n,x,r}(s).
\]
The number is derived from two actual residue frames and precedes
every section. No comparison factor, coefficient formula or matching
frame is an input.
\uses{def:native-tc-weighted-laurent-section-norm,thm:native-tc-product-frame-true-norm,cor:fc-native-laurent-support-frame-values,cor:fc-native-residue-frame-transition}
\end{theorem}
\begin{proof}
The actual proper total-space center supplies a single residue frame
$e_{\mathcal X}$ of the product pullback of $L^{\otimes n}$ in which
all original total-space section norms are their true native coefficient
values. Independently, original properness and invertibility on $X$
choose a genuine local frame around its actual center. Its original
product support frame $e_X$ computes the value of every entire product
section as the weighted maximum of the genuine original coefficient
norms. Both frames are frames of exactly the same original residue
module at the same product Gauss point. Their entire transition
automorphism is multiplication by one actual residue-field unit
$u(e_X,e_{\mathcal X})$. Set $c_{n,x,r}$ equal to its native valuation.
This is strictly positive because the unit has an inverse. The actual
coefficient transition identity now gives the asserted equality for
every original transported section, including zero. No properness or
bounded valuation on the Laurent product is used.
\end{proof}

\begin{corollary}\label{cor:native-tc-weighted-laurent-finite-section-norms}
\lean{EnvelopeProperty.NativeTestConfiguration.TestConfiguration.exists_weightedLaurentPowerFamilyNorm_factor}
\leanok
For any original finite family $G\subset H^0(\mathcal X,M^{\otimes n})$,
let $T_G$ be the actual original finite family of coefficient/exponent
pairs constructed before all points and radii. There is one positive
$c_{n,x,r}$, independent of $G$, for which every such family satisfies
\[
 \max_{s\in G}\lVert s\rVert_{\mathcal X,r,x}
 =c_{n,x,r}\max_{(a,i)\in T_G}\{\lVert a(x)\rVert_{L^{\otimes n}}r^i\}.
\]
The assertion retains empty families and all zero sections.
\uses{thm:native-tc-weighted-laurent-true-section-norms,thm:native-tc-global-laurent-fixed-family}
\end{corollary}
\begin{proof}
Choose the preceding factor before choosing $G$. Taking finite maxima
preserves multiplication by this nonnegative number. The maximum of
the individual coefficient-support maxima is exactly the maximum over
their finite union, which is the defining original family $T_G$.
Finite images and unions do not change a maximum when values repeat.
For the empty family both sides are zero.
\end{proof}

\begin{remark}
The statements compare actual norms in each fixed degree. Compatibility
of the derived factors with tensor powers and the ensuing logarithmic
Fubini--Study and envelope identities require further arguments.
\end{remark}
